\section{Theoretical Background and Related Work}
\label{sec:background}

\subsection{Solar irradiance and forecasting horizons}

Global horizontal irradiance is the total irradiance incident on a horizontal
surface,
\begin{equation}
  \mathrm{GHI}=\mathrm{DHI}+\mathrm{DNI}\cos(\theta_z),
  \label{eq:ghi}
\end{equation}
where DHI and DNI denote diffuse horizontal and direct normal irradiance, and
$\theta_z$ is the solar zenith angle \cite{duffie2013}. The NSRDB provides a
satellite-derived, model-based irradiance reference rather than measurements
collected at the factories \cite{sengupta2018}.

Hourly forecasts preserve the night--day cycle and short atmospheric
variations. Daily aggregation suppresses intraday structure while retaining
weather-driven and seasonal changes. Monthly aggregation emphasizes the
annual cycle but produces very short training series. Accordingly, forecast
horizon must be stated jointly with temporal resolution; terms such as
``short-term'' and ``long-term'' are ambiguous without both quantities.

\subsection{Time-series learning for solar forecasting}

Solar forecasting uses statistical baselines, tree-based learners, recurrent
networks, convolutional models, and attention-based architectures
\cite{voyant2017,voyant2022}. Transparent persistence and seasonal-naive
forecasts remain essential because they reveal whether added architectural
complexity improves upon temporal continuity and seasonality.

TimesNet transforms dominant one-dimensional temporal variations into
two-dimensional representations and processes within-period and between-period
variation with convolutional kernels \cite{wu2023}. This representation is
well suited to signals with multiple periodicities, including the diurnal and
seasonal components of GHI.

Dilated recurrent networks connect a state to selected earlier states through
increasing temporal skips. The shorter computational paths can expose
dependencies at several scales without extending a conventional recurrent
chain one step at a time \cite{chang2017}. Global models additionally share
parameters across related locations while using a site identifier to preserve
location-specific behavior \cite{montero2021}.

\subsection{Industrial transition and climate variability}

\DraftNote{Expand this subsection with verified primary and high-quality
sources on industrial electrification, photovoltaic adoption at manufacturing
sites, the energy transition through 2050/2060, and the documented Ca\c{c}apava
industrial example. Do not state a legal obligation unless the applicable
jurisdiction and regulation are identified.}

Cloud variability deforms otherwise regular solar profiles and generally
increases forecast difficulty. The final literature review will distinguish
irradiance variability, cloud cover, and precipitation instead of treating low
GHI alone as proof of a rainy day.

